\section{The mechanism}\label{sec:mechanism}
The proof follows one line. A sum of squares is a positive semidefinite
Gram matrix on wedge coordinates (\S\ref{sec:mech-wedges}). An exact
identity removes the freedom of that matrix from one fixed combination
of its blocks (\S\ref{sec:mech-identity}).
In a scaling limit, the fixed kernel that remains has a negative direction
(\S\ref{sec:mech-tangent}). Positivity makes the only remaining
Gram-dependent term small (\S\ref{sec:mech-continuation}), and every error is explicit at a finite
scale (\S\ref{sec:mech-estimate}). This section explains each step and
why it works; Sections~\ref{sec:corner}--\ref{sec:finish} give the proofs.

\subsection{Wedges and Pl\"ucker freedom}\label{sec:mech-wedges}
For each symbol label $\lambda$ write $w_\lambda=(x_\lambda,y_\lambda)$ and
\[
 z_{\lambda\lambda'}=w_\lambda\wedge w_{\lambda'}=x_\lambda y_{\lambda'}-x_{\lambda'} y_\lambda .
\]
These wedges vanish simultaneously when the symbol vectors are
proportional. Because $F_N$ vanishes when $x=0$, when $y=0$, and when
$x=y$, every square in a sum of squares vanishes there too, and is
therefore a linear combination of the wedges (Lemma~\ref{lem:frame}).
Hence $F_N$ is SOS exactly when $F_N=z^TQz$ for a positive semidefinite
Gram matrix $Q$ on the wedge space.

Such a $Q$ is far from unique. Wedges of actual symbol vectors obey the
Pl\"ucker identities
\[
 z_{ab}z_{cd}-z_{ac}z_{bd}+z_{ad}z_{bc}=0,
\]
and adding a multiple of one of them changes $Q$ without changing the
polynomial. A negative eigenvalue of one particular Gram matrix therefore
proves nothing: a proof of non-SOS must exclude every Pl\"ucker
modification at once.

This is also where nonnegativity and sums of squares part ways.
Nonnegativity tests $z^TQz$ only on decomposable vectors $z=x\wedge y$,
which satisfy the Pl\"ucker identities; a sum of squares requires some
$Q$ in the Pl\"ucker coset to be positive on all wedge vectors. For large
$N$ we construct a linear functional $\Phi$ on Gram matrices with two
properties: $\Phi(Q)\ge0$ whenever $Q\succeq0$, and $\Phi(Q)<0$ whenever
$Q\succeq0$ represents the corner of $F_N$. No representation can have
both.

\subsection{The corner and the cancellation identity}\label{sec:mech-identity}
The obstruction lives near the two ends of the Toeplitz symbol. The
diagonal at depth $p$ from either end has only $p+1$ nonzero entries,
whatever the ambient order. Diagonals from opposite ends interact in two
disjoint corners of the matrix, and the overlap of two such products is
the minimum of four depths. Retaining the outer $m$ diagonals on each
side therefore gives one polynomial $\calB_m$, the same for every
$N\ge2m$, with product weights $(p+1)(q+1)$ and a minimum kernel
(Lemma~\ref{lem:corner}).

The wedges of the corner are \emph{pure} (both diagonals from the same
end) or \emph{mixed} (one from each end). Averaging a Gram matrix over a
sign change of one side and over the interchange of the two sides loses
no feasible matrix. The averaged matrix has no pure--mixed entries, and
it splits into two pure blocks $P_\pm$, on the sums and differences of
the pure wedges of the two sides, and one mixed block $M$. Each block is
positive semidefinite when the Gram matrix is.

Each block is a fixed part plus a correction. The fixed parts come from
one explicit Gram matrix of $\calB_m$ (Lemma~\ref{lem:corner}). Its two
pure diagonal blocks both equal $D$, a diagonal matrix with the product
weight $2(p+1)(q+1)$ on the pure wedge of depths $p<q$. Its block
coupling the pure wedges of the two sides is $K$, which pairs wedges
with the same gap between their two depths. Its mixed block $B$ consists
of the product weights on the diagonal minus a minimum kernel on mixed
wedges with equal depth totals.

The corrections carry the entire Pl\"ucker freedom. Pl\"ucker relations
among four diagonals from the same end change the pure diagonal blocks
by a \emph{pure four-form} $E$, a matrix whose entries are totally
antisymmetric in the four depths involved. Relations among two
diagonals from each end change the block that couples the pure wedges of
the two sides by a matrix $\Theta$, symmetric after averaging, and
simultaneously change the mixed block by $-\calR\Theta$, where $\calR$
is the realignment described next. The remaining relations couple pure
and mixed wedges, and averaging removes them. Thus
\[
 P_+=D+E+K+\Theta,\qquad P_-=D+E-K-\Theta,\qquad M=B-\calR \Theta.
\]

To compare pure and mixed blocks, identify each matrix $G$ with its
generating kernel $G(\xi,\eta,\zeta,\omega)=f(\xi,\eta)^TGf(\zeta,\omega)$
in four complex variables, where the feature $f$ evaluates a pure wedge
of depths $p<q$ as $\xi^p\eta^q-\xi^q\eta^p$ and a mixed wedge as
$\xi^p\eta^q$. The \emph{realignment} $\calR$ turns a pure kernel $P$ into
a mixed kernel by exchanging the two middle variables,
$(\calR P)(\xi,\eta,\zeta,\omega)=P(\xi,\zeta,\eta,\omega)$. The kernel of a
four-form is alternating in all four variables, so $\calR E=-E$.
Hence both corrections cancel in
\begin{equation}\label{eq:mechanism-cancellation}
 M+\frac{P_++P_-}{2}+\calR P_+=L_m,
\end{equation}
whose right side $L_m=B+D+\calR(D+K)$ is fixed. The identity is exact and
uses no positivity. If the Gram matrix is positive semidefinite, the
first three terms are positive kernels. Realignment does not preserve
positive semidefiniteness, so \eqref{eq:mechanism-cancellation} does not
force $L_m$ to be positive. It reduces the problem to two tasks: find
test points on which $L_m$ is negative, and show that $\calR P_+$ is small
there, uniformly over every allowed Gram matrix.

\subsection{The tangent kernel}\label{sec:mech-tangent}
The test evaluates kernels at nodes $\xi=ie^{-\epsilon a}$, where the
speed $a$ has positive real part. Evaluation at $\xi$ weights depth $p$
by $|\xi|^p=e^{-\epsilon p\Re a}$, so these nodes see
depths of order $1/\epsilon$, deep inside a corner of depth
$m=\epsilon^{-2}$. At this scale the fixed kernel has a limit. After
multiplication by $\epsilon^4$, $L_m$ evaluated at the nodes of the
speed pairs $(a,b)$ and $(a',b')$ tends to
\begin{equation}\label{eq:mechanism-tangent}
 S_F((a,b),(a',b'))=
 2\left(\frac2{A_0^2}+\frac1{C_0^2}-\frac1{A_0C_0}\right),
\end{equation}
where $A_0=(a+\bar a')(b+\bar b')$ and $C_0=(a+b)(\bar a'+\bar b')$.

This tangent kernel has a transparent meaning. For a finite exponential
sum $\phi(s,t)=\sum_\ell c_\ell e^{-a_\ell s-b_\ell t}$ with real coefficients
$c_\ell$, the pairing
$c^*S_Fc$ equals
\begin{equation}\label{eq:energy}
 4\int_{\R_+^2} st|\phi(s,t)|^2\,ds\,dt
 +2\left|\sum_\ell\frac{c_\ell}{(a_\ell+b_\ell)^2}\right|^2
 -2\int_{\R_+^2}|\mathcal H\phi(s,t)|^2\,ds\,dt,
\end{equation}
where $\mathcal H\phi(s,t)=\int_0^\infty\phi(s+u,t+u)\,du$ integrates $\phi$ along
the diagonal ray starting at $(s,t)$. The first two terms are positive.
The third, a diagonal-tail energy, enters with a negative sign. For the
sum given by \eqref{eq:witness}, the pairs $\sigma=\pm1$ combine into
multiples of $e^{-r(s+t)}\sin(j(s-t))$, so $\phi$ is antisymmetric, the
middle term vanishes, and the last term exceeds the first. This is the
negative direction; Lemma~\ref{lem:witness} certifies it by exact
rational arithmetic.

Each exponential alone gives a positive value; only the coherent signed
sum is negative. The test therefore acts on Gram matrices, not on values
of $F_N$, which are nonnegative: it contradicts the positive
semidefiniteness required of every candidate representation.

\subsection{Controlling the realigned term}\label{sec:mech-continuation}
At the test nodes every entry of $\calR P_+$ is a Hermitian diagonal value
$P_+[\xi_\ell,\bar\xi_h]$ of the pure kernel \eqref{eq:realigned-diagonal}. For an
arbitrary positive semidefinite Gram matrix such values can be of order
$\epsilon^{-4}$ at nearby complex nodes (Lemma~\ref{lem:global}); after
the scaling by $\epsilon^4$ this is the size of the main term. Two facts
improve the bound.

First, positivity forces the mixed block to annihilate the complete
antidiagonal sums of low total (Lemma~\ref{lem:means}). This follows by
evaluating all retained variables on a geometric progression and
comparing the first nonzero coefficient of a sum of real squares.
Incomplete high antidiagonals are not forced to vanish; their effect is
exponentially small at the nodes used here (Lemma~\ref{lem:tails}). As
a result, on the conjugate set $\eta=\bar\xi$ the pure kernel is at most
of order $\epsilon^{-2}$, with an explicit finite tail.

Second, a finite positive semidefinite matrix $P_+$ is the Gram matrix of
a holomorphic vector $V$ with $\|V(\xi,\eta)\|^2=P_+[\xi,\eta]$, and
$\log\|V\|^2$ is subharmonic. Two maximum-principle comparisons on
rectangles, one in each speed variable, interpolate between the bounds
$\epsilon^{-4}$ and $\epsilon^{-2}$ and give
\[
 \epsilon^4\|V\|^2\le2^{20}\epsilon^{2^{-17}}
\]
on the region needed by all $1024^2$ entries of the test
(Lemma~\ref{lem:continuation}). The exponent is small because each
comparison retains only a lower bound for the harmonic weight of one
side of a rectangle. The estimate holds for finite matrices of every
size and rank.

\subsection{The finite estimate and the size of the threshold}\label{sec:mech-estimate}
With $m=\epsilon^{-2}$, the finite kernel $\epsilon^4L_m$ differs from
$S_F$ by at most $2^{30}\epsilon$ per entry at the test nodes, and
$\|c\|_1^2<2^{60}$. Pairing \eqref{eq:mechanism-cancellation} with $c$
gives
\[
 0\le c^*\mathsf G_mc
 =\epsilon^4c^*\mathbf L_mc-\epsilon^4c^*\mathbf R_mc
 <-2^{42}+2^{90}\epsilon+2^{80}\epsilon^{2^{-17}},
\]
where $\mathsf G_m$, $\mathbf L_m$ and $\mathbf R_m$ are the matrices of
evaluated entries of $\epsilon^4(M+(P_++P_-)/2)$, $L_m$ and $\calR P_+$;
the first is positive semidefinite. This is Theorem~\ref{thm:finite-scale}.

The threshold is large for a quantitative reason. The integer
coefficients make the main term look large, $c^*S_Fc<-2^{42}$, but every
entrywise error is multiplied by $\|c\|_1^2$, and relative to $\|c\|_1^2$
the certified negativity is small: $-c^*S_Fc/\|c\|_1^2$ lies between
$2^{-18}$ and $2^{-16}$, on the scale of $10^{-5}$. The error
$2^{90}\epsilon$ is harmless, but $2^{80}\epsilon^{2^{-17}}$ decays
extremely slowly and falls below $2^{42}$ only when
$\epsilon<2^{-38\cdot2^{17}}$. The largest dyadic $\epsilon$ with a
negative budget is $2^{-\ScaleExponent}$, which gives
$m=\epsilon^{-2}=\CornerDepth$ (Appendix~\ref{app:constants}). These
constants explain the sufficient threshold; they do not predict the
location of the first non-SOS order.

\subsection{Why small orders are sums of squares}\label{sec:mech-small}
At a fixed small order the Pl\"ucker freedom can make a Gram matrix
positive on the complement of its forced kernel. The exact certificates
through order $50$ do precisely this, with strictly positive rational
bounds on the relevant compressed blocks (Section~\ref{sec:positive}).
The analytic obstruction requires a deep corner and a geometric decay
length $1/\epsilon$ that fits within it, at a scale where the errors
above fall below the main term; no estimate here makes it effective at
small orders. Positive certificates must nevertheless
lose their relative margin as the order grows
(Section~\ref{sec:exchange}), and the natural induction through positive
increments fails at all large orders (Section~\ref{sec:boundary}).
